

\subsection{Design Weaknesses}

If you are looking for reasons that \UT{} might not work, or for the most valuable contributions, let's cut to the chase.

\textbf{Dapp-chains \& PoS:} \autoref{sec:dapp-chains}.
Dapp-chains are based on the assumption that a secure PoS method exists when adapted for one-way PoR.
No specific PoS method is provided and it's not clear that PoS + PoR actually works without a clear design for such a system.
The other claims about dapp-chains and \UT{2+} scaling are contingent on a suitable PoS method being found or created.
PoS is insecure when used in isolation and it isn't clear that PoS can be secure, even with the help of PoR and a purely PoW simplex.
Even if a suitable method of PoS is found, PoS + PoR is inefficient compared to PoW + PoR and capacity estimates are overly optimistic.

\textbf{\UTinf{1} Block Availability \& $O(c^2 \log n)$ Bandwidth:} \autoref{sec:tiling-bandwidth}.
Without a solution to block availability, \UTinf{} doesn't work for $d \gtrsim 4$.
It works to expand capacity by up to $\sim 10\times$ but not beyond.
Even that might not work because the security model is untested and there are no comparable PoW-based real-world examples to examine.
$O(c^2 \log n)$ bandwidth for validating nodes is too high.
$O(n c^{-i})$ bandwidth for root-tile miners is too high.

\textbf{A Key Point for \UT1 Is Flawed:}
Several key components are required for \UT1 to work.
There are \ul{no known decisive criticisms} but maybe you can find one?
PoR and Conversion of Work (\autoref{sec:proof-of-reflection}), Simplex Security (\autoref{sec:simplex-security}), and Intra-Simplex Cross-Chain Transactions (\autoref{sec:spv-in-ut}) are three such examples.

% \subsection{} answered section?

\textbf{\UT1 $O(c^2)$ Bandwidth:} \autoref{sec:availability-of-blocks}.
Depends on chosen $k$ (block size) and \ul{is mitigated} with a suitable $k$.
A $k$ of 3000 B/s has plenty of breathing room, and 20,000 B/s is on the high end of what could possibly work.
For comparison, Bitcoin has $k \approx 1700$.


\priorityTODO[h]{list weaknesses and current problems}

% \subsubsection{Not Weaknesses}

% \textbf{Cryptocurrency Isn't Real Money:} This is true in the sense that all cryptocurrency is a type of fiat, but then government money isn't real money either.
% Regardless, the value of cryptocurrency is in its ability to \emph{act like money} and the \emph{utility} it provides.
% So long as people are using cryptocurrency with rational judgement to preserve or create value and increase profits, it will maintain that utility to some degree.
% Gold will always be a better form of money in an absolute sense, but, in the \emph{context of the internet}, PoW cryptocurrency provides something that gold cannot: \emph{freedom from tyranny} (among other things, too, of course).

% \textbf{PoW Using Electricity Is Bad:} This does not explain why UT will fail, and it is not a weakness unless you are an environmentalist (which I am not).
% \ul{Mitigated} by: prioritizing one's own life and flourishing over nature; understanding the connections between energy use, standard of living, supply and demand, prices, capitalism generally, and scarcity's transformative nature under capitalism; recognizing that increasing energy demand is a necessary part of ongoing and accelerating scientific and technological progress and therefore cannot be intrinsically bad; realizing that the whole world is made of natural resources; valuing highly accessible, metaphysically secure non-government money and the important role it will play in safeguarding one's life and future; using a philosophically optimistic epistemology; and learning that atoms do not have a purpose until a thinking being gives them one, especially through the conscious act of transformation in the fulfillment of their \ul{own} \emph{values}.

% \textbf{Cryptocurrency Is Full of Scams:} Yes. Those are bad.
% Judging a technology by a subset of its users is also bad, though.
% It is the underlying dishonesty and fraud that is the problem, not the technology itself.
% That does not mean we should ignore the problem, but here is not the right place to address it.
